% ===== PRÉAMBULE CANONIQUE — à coller VERBATIM en tête de CHAQUE .tex =====
\documentclass[11pt,a4paper]{article}
\usepackage[utf8]{inputenc}\usepackage[T1]{fontenc}\usepackage{lmodern}
\usepackage[french]{babel}
\usepackage{amsmath,amssymb,mathtools}\usepackage{icomma}
\usepackage{siunitx}\sisetup{locale=FR}  % \qty{}{}, \si{}, \ang{}
\usepackage{xcolor}\usepackage{colortbl}
\usepackage{tikz}\usepackage{pgfplots}\pgfplotsset{compat=1.18}
\usepackage[siunitx,europeanresistors]{circuitikz} % circuits (élec.)
\usepackage[most]{tcolorbox}\usepackage{enumitem}\usepackage{array,booktabs}
\usepackage[a4paper,margin=2.2cm]{geometry}
\usetikzlibrary{arrows.meta,calc,angles,quotes,positioning,patterns,%
  backgrounds,shapes.geometric,decorations.pathmorphing,decorations.markings,intersections}
\definecolor{primary}{RGB}{222,49,99}\definecolor{primarydark}{RGB}{150,22,55}
\definecolor{defcol}{RGB}{0,114,178}\definecolor{methcol}{RGB}{52,92,148}
\definecolor{excol}{RGB}{0,135,90}\definecolor{attncol}{RGB}{200,40,55}
\tcbset{boxbase/.style={enhanced,breakable,boxrule=0.9pt,arc=2.5pt,
  left=3mm,right=3mm,top=2.3mm,bottom=2.3mm,fonttitle=\bfseries,coltitle=white}}
% --- boîtes TUTORIEL (noms EXACTS, ne pas renommer) ---
\newtcolorbox{cle}[1][]{boxbase,colback=defcol!6,colframe=defcol,
  title={\textsc{À retenir}\ifx\relax#1\relax\relax\else\textmd{ : \itshape #1}\fi}}
\newtcolorbox{methode}[1][]{boxbase,colback=methcol!7,colframe=methcol,sharp corners,
  title={\textsc{Méthode}\ifx\relax#1\relax\relax\else\textmd{ --- \itshape #1}\fi}}
\newtcolorbox{exemplebox}[1][]{boxbase,colback=excol!5,colframe=excol,
  title={\textsc{Exemple}\ifx\relax#1\relax\relax\else\textmd{ : \itshape #1}\fi}}
\newtcolorbox{piege}[1][]{boxbase,colback=attncol!6,colframe=attncol,
  title={\textsc{Piège}\ifx\relax#1\relax\relax\else\textmd{ : \itshape #1}\fi}}
\newtcolorbox{remarque}[1][]{boxbase,colback=black!4,colframe=black!45,
  title={\textsc{Remarque}\ifx\relax#1\relax\relax\else\textmd{ : \itshape #1}\fi}}
% --- boîtes EXERCICE / CORRIGÉ (noms EXACTS, auto-comptées) ---
\newtcolorbox[auto counter]{exo}[1][]{boxbase,colback=primary!4,colframe=primary,
  title={\textsc{Exercice}~\thetcbcounter\ifx\relax#1\relax\relax\else\textmd{ --- \itshape #1}\fi}}
\newtcolorbox[auto counter]{corrige}[1][]{boxbase,colback=excol!4,colframe=excol,
  title={\textsc{Corrigé}~\thetcbcounter\ifx\relax#1\relax\relax\else\textmd{ --- \itshape #1}\fi}}
\newcommand{\vv}[1]{\vec{#1}}\newcommand{\ds}{\displaystyle}
\newcommand{\R}{\mathbb{R}}\newcommand{\N}{\mathbb{N}}\newcommand{\Z}{\mathbb{Z}}
% (un \usepackage{fancyhdr}+en-tête est possible ; il est ignoré côté web)
% =========================================================================
\begin{document}

\begin{center}
{\large\bfseries\color{excol} Thème 2 — Corrigé Difficile}\\[-2pt]
{\color{excol}\rule{5cm}{1pt}}
\end{center}

\begin{corrige}[dans quel sens le point P se déplace-t-il ?]
\begin{enumerate}[label=\alph*)]
  \item Le point $P$ va \textbf{monter}.
  \item L'onde vient de la \emph{gauche} (elle avance vers la droite). Le voisin de $P$ situé à
        gauche est plus haut (il est plus proche du sommet de la bosse). À l'instant suivant,
        $P$ prend la valeur de ce voisin : il monte donc. Autrement dit, la bosse \emph{arrive}
        sur $P$, qui commence par s'élever. (Un point du flanc \emph{arrière}, lui, descendrait.)
\end{enumerate}
\end{corrige}

\begin{corrige}[d'un décalage spatial au déphasage]
Un décalage spatial $d$ correspond à un déphasage $\Delta\varphi=2\pi\,d/\lambda$.
\begin{enumerate}[label=\alph*)]
  \item Décalage $d=\lambda/4$ :
        $\Delta\varphi=2\pi\times\dfrac{\lambda/4}{\lambda}=2\pi\times\dfrac14=\dfrac{\pi}{2}\ \si{\radian}.$
  \item L'opposition de phase correspond à $\Delta\varphi=\pi$, soit $d=\lambda/2=\SI{4}{\metre}$.
\end{enumerate}
\end{corrige}

\begin{corrige}[de l'enregistrement à l'équation, puis aux vitesses]
Le graphe est en fonction du temps : on y lit $A$ et $T$. Puis $\omega=2\pi/T$, $k=2\pi/\lambda$.
\begin{enumerate}[label=\alph*)]
  \item $A=\SI{4}{\centi\metre}=\SI{0.04}{\metre}$ ; $T=\SI{0.25}{\second}$. Donc
        \[ \omega=\frac{2\pi}{T}=\frac{2\pi}{0{,}25}=8\pi\approx\SI{25.1}{\radian\per\second},
           \quad
           k=\frac{2\pi}{\lambda}=\frac{2\pi}{0{,}5}=4\pi\approx\SI{12.6}{\radian\per\metre}. \]
  \item $y(x,t)=0{,}04\,\sin\!\left(8\pi\,t - 4\pi\,x\right)$ (SI).
  \item $v_{\max}=A\omega=0{,}04\times 8\pi=0{,}32\pi\approx\SI{1.0}{\metre\per\second}$ ;
        \[ a_{\max}=A\omega^{2}=0{,}04\times(8\pi)^{2}=0{,}04\times 64\pi^{2}
           \approx\SI{25.3}{\metre\per\second\squared}. \]
\end{enumerate}
\end{corrige}

\begin{corrige}[déphasage supérieur à un tour]
On calcule le déphasage brut, puis on retire des tours entiers ($2\pi$) tant que possible.
\begin{enumerate}[label=\alph*)]
  \item $\Delta\varphi=2\pi\,\dfrac{d}{\lambda}=2\pi\times\dfrac{10}{4}=2\pi\times 2{,}5=5\pi\ \si{\radian}.$
  \item $5\pi = 4\pi + \pi = 2\times(2\pi) + \pi$ : on retire deux tours complets, il reste
        $\Delta\varphi\equiv\pi$. C'est donc une \textbf{opposition de phase}
        (cohérent avec $d=\SI{10}{\metre}=2{,}5\,\lambda=(2\times2+1)\lambda/2$).
\end{enumerate}
\end{corrige}

\begin{corrige}[le point peut-il dépasser l'onde ?]
\begin{enumerate}[label=\alph*)]
  \item $v=\lambda f=0{,}2\times 50=\SI{10}{\metre\per\second}$ ;
        $\omega=2\pi f=100\pi$, donc
        $v_{\max}=A\omega=0{,}02\times 100\pi=2\pi\approx\SI{6.3}{\metre\per\second}.$
  \item La célérité $v=\SI{10}{\metre\per\second}$ est \textbf{plus grande} que
        $v_{\max}\approx\SI{6.3}{\metre\per\second}$ : ici l'onde se propage plus vite que le
        point n'oscille.
  \item $v_{\max}=v \iff A\omega=\lambda f \iff A\,(2\pi f)=\lambda f \iff A=\dfrac{\lambda}{2\pi}.$
        Numériquement $A=\dfrac{0{,}2}{2\pi}\approx\SI{0.032}{\metre}\approx\SI{3.2}{\centi\metre}$ :
        au-delà de cette amplitude, le point du milieu irait, à son maximum, plus vite que l'onde.
\end{enumerate}
\end{corrige}

\end{document}
